\documentclass[10pt,a4paper]{amsart}
\usepackage[margin=19mm]{geometry}
\usepackage{amsmath,amssymb,mathtools}
\usepackage[scr=rsfs,cal=euler]{mathalfa}
\usepackage{xfrac}
\usepackage[unicode,hidelinks,bookmarks=false]{hyperref}
\newcommand{\ie}{i.e.\ }
\newcommand{\cf}{cf.\ }
\newcommand{\resp}{respectively\ }
\newcommand{\Subsection}{\subsection}
\providecommand{\nobreakdash}{\nobreak\hbox{-}}
\setlength{\emergencystretch}{2em}
\multlinegap0pt
\usepackage{bm}
\usepackage{bbm}
\usepackage{dsfont}
\usepackage{xspace}
\usepackage{cancel}
\usepackage{mathtools}
\usepackage{amsthm}
\makeatletter
\@ifundefined{theorem}{\newtheorem{theorem}{Theorem}}{}
\makeatother
\title[Rolls and Snaking in a Swift-Hohenberg Equation with Non-smo]{Rolls and Snaking in a Swift-Hohenberg Equation with Non-smooth Nonlinearity}
\author{L\"utfiye Masur, Jens D. M. Rademacher}
\date{Source: arXiv:2609.24355v1 (CC BY 4.0)}
\begin{document}
\maketitle

\noindent\emph{Source and licence.} Rolls and Snaking in a Swift-Hohenberg Equation with Non-smooth Nonlinearity. Version arXiv:2609.24355v1 (CC BY 4.0), distributed under the Creative Commons Attribution 4.0 International licence, as recorded in the arXiv licence metadata for this exact version and shown on the abstract page https://arxiv.org/abs/2609.24355v1. Full rights notice: licenses/arxiv-2609.24355-CC-BY-4.0.txt.\par\smallskip

\noindent\textit{Editorial scope.} Counted unit: the Theorem whose source label is \texttt{theorem: reduced system nu>1/2} in the article (source0.tex lines 731--740, source label \texttt{theorem: reduced system nu>1/2}), delivered together with the article's own complete original proof of it (source0.tex lines 742--881, one \texttt{proof} environment of 140 lines). No mathematics is written, repaired or re-derived: statement and proof are the article's own text line for line. Required context retained uncounted: SH1 (lines 644--646); neumann conditions (lines 672--674); matching condition (lines 676--678); smoothness condition (lines 680--682); eq:u+-sol (lines 685--687); eq: gamma's (lines 715--728). Every definition, display and label of those lines is retained unchanged. Omitted from this package and not redistributed: 1--643, 647--671, 675--675, 679--679, 683--684, 688--714, 729--730, 741--741, 882--1242. Nothing inside a retained excerpt is omitted or altered: every retained body is delivered exactly as the article's lines are written, so no omission receipt of any kind accompanies this package. Citations: the retained excerpts carry no citation command at all, so no bibliography entry of the article is transcribed and no reference is redistributed. Reference adaptation: none was needed and none was made; every reference inside the delivered unit resolves inside it, so no unresolved reference or citation is produced. Printed slips kept literal: no printed word, symbol, hypothesis or number of the counted statement or of its proof was altered. Layout and class: the article's own class is \texttt{article} with packages including babel, geometry, amsmath, amssymb, amsthm, graphicx, hyperref, xcolor, subfig, todonotes, authblk; the wrapper is the lane's pinned \texttt{amsart} editorial wrapper at 10pt on A4, which restates the theorem-like environments the retained text needs and adds the article's own macro definitions that the retained text needs (none), with extra wrapper packages (bm, bbm, dsfont, xspace, cancel, mathtools). The delivered numbering is the unit's own. Assets: no retained span contains a figure environment, a graphics-inclusion command, a PDF-page inclusion, a table or a TikZ picture; no image file is copied into this package, the package declares no asset dependency and no image file is redistributed with it. Licence: version arXiv:2609.24355v1 of this article is distributed under the Creative Commons Attribution 4.0 International licence, as recorded in the arXiv licence metadata for this exact version and shown on the abstract page https://arxiv.org/abs/2609.24355v1. A full rights notice is delivered at licenses/arxiv-2609.24355-CC-BY-4.0.txt. Characters: every retained excerpt is pure ASCII, so the delivered engine, XeLaTeX with its default Latin Modern text font, reports no missing character.
		\begin{equation} \label{SH1}
			0 = - (1+\partial_x^2)^2 u + \mu u + \nu |u|, 
		\end{equation}

		\begin{equation}\label{neumann conditions}
			u_{-}^{(i)} (-L_{-})=0 \text{ and } u_{+}^{(i)} (L_{+})=0 \text{ for } i=1,3,
		\end{equation}

		\begin{equation}\label{matching condition}
			u_{+} (0)=u_{-} (0) =0
		\end{equation}

		\begin{equation}\label{smoothness condition}
			u_{+}^{(j)} (0)=u_{-}^{(j)} (0) \text{ for } j=1,2,3 .           	
		\end{equation}

		\begin{equation}\label{eq:u+-sol}
			u_\pm(x) = a_{\pm 1} e^{i \kappa_{\pm 1}x} +  a_{\pm 2} e^{-i \kappa_{\pm 1}x} +  a_{\pm 3} e^{i \kappa_{\pm 2}x} +  a_{\pm 4} e^{-i \kappa_{\pm 2}x}
		\end{equation}

	\begin{align} \label{eq: gamma's}
		\begin{split}
		\eta = & \mathrm{Re}(i \kappa_{-2}), \quad \beta = \mathrm{Im}(i \kappa_{-2}),\\
			\gamma_{31} =& (-1 + e^{2 \eta (L_+ -\pi)}) (\kappa_{+1}^2 - \kappa_{+2}^2)-2 \eta (1 + e^{2 \eta (L_+ -\pi)}) \kappa_{+1} \tan(\kappa_{+1} L_+) \\
			&  +2 \eta (1 + e^{2 \eta (L_+ -\pi)}) \kappa_{+2} \tan(\kappa_{+2} L_+), \\
			\gamma_{32} =& \eta (1 + e^{2 \eta (L_+ -\pi)}) (\kappa_{+1}^2 - \kappa_{+2}^2) + (-1 + e^{2 \eta (L_+ -\pi)}) \kappa_{+1} (\eta^2 - \beta^2 + \kappa_{+1}^2) \tan(\kappa_{+1} L_+) \\
			& - (-1 + e^{2 \eta (L_+ -\pi)}) \kappa_{+2} (\eta^2 - \beta^2 + \kappa_{+2}^2) \tan(\kappa_{+2} L_+), \\
			\gamma_{41} =& (3 \eta^2 - \beta^2)(-1 + e^{2 \eta (L_+ -\pi)})(\kappa_{+1}^2 - \kappa_{+2}^2) +2 \eta (1 + e^{2 \eta (L_+ -\pi)}) \kappa_{+1}^3 \tan(\kappa_{+1} L_+) \\
			& -2 \eta (1 + e^{2 \eta (L_+ -\pi)}) \kappa_{+2}^3 \tan(\kappa_{+2} L_+), \\
			\gamma_{42} =& \eta (\eta^2 -3\beta^2) (1 + e^{2 \eta (L_+ -\pi)}) (\kappa_{+1}^2 - \kappa_{+2}^2) \\
			& + (-1 + e^{2 \eta (L_+ -\pi)}) \kappa_{+1} \big( \eta^4 + \beta^4 -\beta^2 \kappa_{+1}^2 + \eta^2 (2\beta^2 + \kappa_{+1}^2) \big) \tan(\kappa_{+1} L_+) \\
			& - (-1 + e^{2 \eta (L_+ -\pi)}) \kappa_{+2} \big( \eta^4 + \beta^4 -\beta^2 \kappa_{+2}^2 + \eta^2 (2\beta^2 + \kappa_{+2}^2) \big) \tan(\kappa_{+2} L_+).
		\end{split}
	\end{align} 

	\begin{theorem}\label{theorem: reduced system nu>1/2}
For given $\nu > 1/2$, the reflection symmetric $2\pi$-periodic real valued solutions of \eqref{SH1} with two sign changes over one period with $\mu\in (-\nu,1-\nu)$  are in 1-to-1 correspondence with solutions $(\mu,L_+)\in (-\nu,1-\nu)\times [0,\pi]$ to 
		\begin{align}
			\begin{split}\label{reduced system}
				&f_\nu (\mu,L_+) := \gamma_{31} \gamma_{42} - \gamma_{41} \gamma_{32} = 0 \\
				&g_\nu (\mu,L_+) := \beta \cos\big(\beta (L_+ -\pi)\big) \gamma_{31} + \sin\big(\beta (L_+ - \pi)\big) \gamma_{32} = 0
			\end{split}
		\end{align}
		with $\gamma_{ij}, \beta, \eta$ defined in \eqref{eq: gamma's}.
\end{theorem}

		\begin{proof} 
		For $-\nu < \mu < 1-\nu$ and $\nu>1/2$ the solution $u_+$ can be written as
		\begin{align}
			\begin{split}
				u_+(x) =& (a_{+1} + a_{+2}) \cos (\kappa_{+1} x) + i ( a_{+1} - a_{+2} ) \sin (\kappa_{+1} x) \\
				&+ ( a_{+3} + a_{+4}) \cos (\kappa_{+2} x) + i ( a_{+3} - a_{+4}) \sin (\kappa_{+2} x) .
			\end{split}
		\end{align}
		Since we search for real solutions, $a_{+1} + a_{+2}$, $ a_{+3} + a_{+4}$ are real and $a_{+1} - a_{+2}$, $a_{+3} - a_{+4}$ are purely imaginary so that $\mathrm{Im}(a_{+1}) = -\mathrm{Im}(a_{+2})$, $\mathrm{Re}(a_{+1}) = \mathrm{Re}(a_{+2})$ and, $\mathrm{Im}(a_{+3}) = -\mathrm{Im}(a_{+4})$, $\mathrm{Re}(a_{+3}) = \mathrm{Re}(a_{+4})$. By the symmetry relations of the imaginary parts of $a_{\pm j}'s$, the matching condition in \eqref{matching condition} gives
		\begin{equation}
			u_+(0) = 2 \mathrm{Re} (a_{+1}) + 2 \mathrm{Re} (a_{+3}) = 0, 
		\end{equation}
		so that $\mathrm{Re}(a_{+1}) = -\mathrm{Re}(a_{+3})$.
		The analogous calculations can be done for $u_{-}(x)$  
		so that 
		\begin{equation}
			\begin{tabular}{p{0.25\linewidth} p{0.25\linewidth}}
				\parbox{\linewidth}{$a_{+1} = -i r_{+1} - s_+ $ \\ $a_{+2} = i r_{+1} - s_+ $ \\ $a_{+3} = -i r_{+2} + s_+ $ \\ $a_{+4} = i r_{+2} + s_+ $}&
				\parbox{\linewidth}{$a_{-1} = -i r_{-1} - s_- $ \\ $a_{-2} = i r_{-2} + s_- $ \\ $a_{-3} = -i r_{-2} + s_- $ \\ $a_{-4} = i r_{-1} - s_- $}
			\end{tabular}
		\end{equation}
		where $r_{\pm j}$ and $s_{\pm}$ are real numbers. Solutions \eqref{eq:u+-sol}  
		are thus of the form
		\begin{align}\label{new u's}
			\begin{split}
				u_+(x) =& 2 r_{+1} \sin(\kappa_{+1}x) + 2 r_{+2} \sin(\kappa_{+2}x)  -2 s_+ \cos(\kappa_{+1}x) + 2 s_+ \cos(\kappa_{+2}x)\\
				u_-(x) =& 2 e^{-\eta x} r_{-1} \sin(\beta x) + 2 e^{\eta x} r_{-2} \sin(\beta x) -2 e^{-\eta x} s_- \cos(\beta x) + 2 e^{\eta x} s_- \cos(\beta x),
			\end{split}
		\end{align} 
		where $i \kappa_{-1} = -\eta + i \beta$ and $i \kappa_{-2} = \eta + i \beta$. Notably, the matching conditions \eqref{matching condition} are satisfied. (For $\nu\geq 1/2$ this form is valid only in the parameter regime $-\nu < \mu < 1-\nu$.) 
		 
		 We are left with the 8 unknowns $L_+$, $\mu$ and $r_{\pm j}$, $s_{\pm}$, $1 \leq j \leq 2$, for the six remaining equations  \eqref{neumann conditions}, \eqref{smoothness condition}. It is convenient to write the latter in the form  
		\begin{equation}\label{system}
			\begin{pmatrix}
				\tilde{u}'_+(L_+) & \begin{matrix} 0 & 0 & 0\end{matrix} \\
				\tilde{u}'''_+(L_+) & \begin{matrix} 0 & 0 & 0\end{matrix} \\
				\begin{matrix} 0 & 0 & 0\end{matrix} & \tilde{u}'_-(-L_-) \\
				\begin{matrix} 0 & 0 & 0\end{matrix} & \tilde{u}'''_-(-L_-) \\
				\tilde{u}'_+(0) & -\tilde{u}'_-(0) \\
				\tilde{u}''_+(0) & -\tilde{u}''_-(0) \\
				\tilde{u}'''_+(0) & -\tilde{u}'''_-(0) 
			\end{pmatrix}
			\begin{pmatrix}
				r_{+1} \\ 
				r_{+2} \\
				s_+ \\
				r_{-1} \\
				r_{-2} \\
				s_-
			\end{pmatrix} = 0
		\end{equation} 
		with $1\times 3$ row vectors $u_+^{(k)}(0)$, $k=1,2,3$, resulting from substituting \eqref{new u's} into \eqref{neumann conditions}, \eqref{smoothness condition} (and removing common constant factors). One readily computes these vectors as 
		\begin{equation}
			\begin{tabular}{p{0.31\linewidth} p{0.5\linewidth}}
				\parbox{\linewidth}{$\tilde{u}'_+(0) = \begin{pmatrix}
						\kappa_{+1}, & \kappa_{+2}, & 0
					\end{pmatrix}$ \\
					$\tilde{u}''_+(0) = \begin{pmatrix}
						0, & 0, & \kappa_{+1}^2 - \kappa_{+2}^2  
					\end{pmatrix}$ \\
					$\tilde{u}'''_+(0) = \begin{pmatrix}
						-\kappa_{+1}^3, & -\kappa_{+2}^3, & 0 
					\end{pmatrix}$
				}&
				\parbox{\linewidth}{$\tilde{u}'_-(0) = \begin{pmatrix}
						\beta, & \beta, & 2\eta
					\end{pmatrix}$ \\
					$\tilde{u}''_-(0) = \begin{pmatrix}
						-2 \eta \beta, & 2 \eta \beta, & 0
					\end{pmatrix}$ \\
					$\tilde{u}'''_-(0) = \begin{pmatrix}
						3 \eta^2 \beta - \beta^3, & 3 \eta^2 \beta - \beta^3, & 2 \eta^3 -6 \eta \beta^2
					\end{pmatrix}$
				}	
			\end{tabular}
		\end{equation}
		as well as  
		\begin{align}
			\begin{split}
				\tilde{u}'_+(L_+) =& \begin{pmatrix}
					\kappa_{+1} \cos(\kappa_{+1} L_+), & \kappa_{+2} \cos(\kappa_{+2} L_+), & \kappa_{+1} \sin(\kappa_{+1} L_+) - \kappa_{+2} \sin(\kappa_{+2} L_+)
				\end{pmatrix} \\
				\tilde{u}'''_+(L_+) =& \begin{pmatrix}
					-\kappa_{+1}^3 \cos(\kappa_{+1} L_+), & -\kappa_{+2}^3 \cos(\kappa_{+2} L_+), & -\kappa_{+1}^3 \sin(\kappa_{+1} L_+) + \kappa_{+2}^3 \sin(\kappa_{+2} L_+)
				\end{pmatrix} \\
			\end{split}
		\end{align}
		and, using $L_- = \pi - L_+$, and abbreviating $c_+:=\cos (\beta(L_+ - \pi)), s_+:=\sin (\beta(L_+ - \pi))$, $e_\pm:=e^{\pm\eta (L_+ - \pi)}$, 
		\begin{align}
			\begin{split}
				\tilde{u}'_-(-L_-) =&\left(\beta e_- c_+-\eta e_- s_+, 
				 \beta e_+ c_++\eta e_+ s_+, 
				 (\eta e_-  +\eta e_+) c_+ + (\beta e_-  -\beta e_+) s_+\right) \\
				\tilde{u}'''_-(-L_-) =& \big(
				(3\eta^2 \beta e_-  - \beta^3 e_-) c_+ - ( \eta^3 e_-  - 3\eta \beta^2 e_-) s_+, 
				(3\eta^2 \beta e_+  -\beta^3 e_+) c_+ + (\eta^3 e_+ - 3\eta \beta^2 e_+) s_+,\\
				&
				(\eta^3 e_-  - 3\eta \beta^2 e_-  + \eta^3 e_+  - 3\eta \beta^2 e_+) c_+ 
				+ (3\eta^2 \beta e_-  - \beta^3 e_-  - 3\eta^2 \beta e_+  + \beta^3 e_+ ) s_+				
				\big).
			\end{split}
		\end{align}
		Solving first two and last three rows in \eqref{system} yields
		\begin{equation} \label{eq:rs}
		\begin{aligned}
			r_{+1} &= 
			\frac{
					\kappa_{+2} (1-\kappa_{+2}^2) r_{+2} \cos(\kappa_{+2} L_+) 
					+ \kappa_{+1} (1 - \kappa_{+1}^2) s_+ \sin(\kappa_{+1} L_+)  
					+ \kappa_{+2} (\kappa_{+2}^2-1) s_+ \sin(\kappa_{+2} L_+) 
			}{
				\kappa_{+1} (-1 + \kappa_{+1}^2) \cos(\kappa_{+1} L_+)
			},\\
			r_{+2} &= \frac{s_+ \sin(\kappa_{+2} L_+)}{\cos(\kappa_{+2} L_+)},\\
			r_{-1} &= -\frac{\kappa_{+1}^2 - \kappa_{+2}^2}{2 \eta \beta} s_+ + r_{-2},\\
			r_{-2} &= -
			\frac{
					4 \eta^2 s_- -\kappa_{+1}^2 s_+ + \kappa_{+2}^2 s_+ 
					+ 2\eta \kappa_{+1} s_+ \tan(\kappa_{+1} L_+) 
					- 2\eta \kappa_{+2} s_+ \tan(\kappa_{+2} L_+)
			}{4 \eta \beta},\\
			s_- &= -
			\frac{
				\begin{aligned}
					\kappa_{+1}(3\eta^2 -\beta^2 +\kappa_{+1}^2) s_+ \tan(\kappa_{+1} L_+) 
					-\kappa_{+2}(3\eta^2 -\beta^2 +\kappa_{+2}^2) s_+ \tan(\kappa_{+2} L_+)
				\end{aligned}				
			}{4\eta (\eta^2 + \beta^2)}.
		\end{aligned}\end{equation}		
The remaining third and fourth rows in \eqref{system} thus read
		\begin{equation}\label{eq:fs}
		\begin{aligned}
			\frac{e^{\eta (-L_+ + \pi)}s_+}{4 \eta \beta} \bigl( \beta \cos(\beta(L_+-\pi)) \gamma_{31} + \sin(\beta(L_+-\pi)) \gamma_{32} \bigl) =0\\
			\frac{e^{\eta (-L_+ + \pi)}s_+}{4 \eta \beta} \bigl( \beta \cos(\beta(L_+-\pi)) \gamma_{41} + \sin(\beta(L_+-\pi)) \gamma_{42} \bigl) =0
		\end{aligned}
		\end{equation}
		where $\eta$, $\beta$ and $e^{\eta(-L_+ + \pi)}$ are non-zero. 
		Hence, \eqref{eq:fs} equivalently reduces to \eqref{reduced system} as claimed. 
		If $s_+ = 0$, then $s_{-}$ and all $r_{\pm j}$ vanish, which implies $u(x)=0$ so that no nontrivial periodic solution exists in this case.
	\end{proof}

\end{document}
